% ============================================================
%  code_tta_all.py 구조 설명 슬라이드
%  저장 위치: manuscript/slides/code_tta_all.tex
%  컴파일: pdflatex code_tta_all.tex  (twice)
% ============================================================
\documentclass[aspectratio=169]{beamer}
\usetheme{Madrid}
\usecolortheme{seahorse}

\usepackage{booktabs}
\usepackage{amsmath}
\usepackage{amssymb}
\usepackage{xcolor}
\usepackage{array}
\usepackage{tikz}
\usepackage{pgfplots}
\pgfplotsset{compat=1.18}
\usetikzlibrary{arrows.meta, positioning, shapes.geometric, fit,
                decorations.pathreplacing, calc, backgrounds}
\usepackage{fontspec}
\usepackage{xeCJK}
\setCJKmainfont[AutoFakeSlant=0.2, AutoFakeBold=2]{Noto Sans CJK KR}
\setCJKsansfont[AutoFakeSlant=0.2, AutoFakeBold=2]{Noto Sans CJK KR}
\setCJKmonofont{Noto Sans CJK KR}

% ── 색상 정의 ────────────────────────────────────────────────
\definecolor{expblue}{RGB}{70,130,200}
\definecolor{exporange}{RGB}{220,140,40}
\definecolor{expgreen}{RGB}{60,160,80}
\definecolor{exppurple}{RGB}{160,80,200}
\definecolor{lightgray}{RGB}{240,240,240}
\definecolor{attred}{RGB}{200,60,60}
\definecolor{benignblue}{RGB}{60,100,200}
\definecolor{jsdcolor}{RGB}{180,60,180}

\setbeamertemplate{navigation symbols}{}
\setbeamerfont{footnote}{size=\tiny}

\title{\texttt{code\_tta\_all.py} 구조 분석}
\subtitle{Inclusive Expert + MATI/SADE 기반 TTA 설계}
\author{westpark}
\date{2026-06-05}
\institute{CIC-IDS2017 \quad (2.83M samples, 15 classes)}

% ============================================================
\begin{document}

\begin{frame}
  \titlepage
\end{frame}

% ── 목차 ─────────────────────────────────────────────────────
\begin{frame}{목차}
  \tableofcontents
\end{frame}

% ============================================================
\section{전체 파이프라인}
% ============================================================

\begin{frame}{전체 파이프라인 개요}
\begin{center}
\begin{tikzpicture}[
  node distance=0.55cm and 0.7cm,
  box/.style={draw, rounded corners=3pt, minimum width=2.2cm,
              minimum height=0.7cm, font=\small, align=center,
              fill=lightgray},
  arr/.style={-{Stealth[scale=0.8]}, thick},
  phase/.style={draw=none, font=\scriptsize\bfseries, align=center}
]

% ─ Phase 1: 데이터 ─
\node[box, fill=blue!10] (data) {원시 데이터\\{\scriptsize CIC-IDS2017}};

\node[box, fill=blue!10, right=of data] (split)
  {Train/Val/Test 분리\\{\scriptsize 60/20/20, stratified}};

% ─ Phase 2: 학습 ─
\node[box, fill=green!15, right=of split] (caps)
  {Inclusive Caps\\계산\\{\scriptsize $k$개 경계값}};

\node[box, fill=green!15, right=of caps] (train)
  {Expert 학습\\{\scriptsize XGBoost $\times k$}\\{\scriptsize val도 동일 cap}};

\node[box, fill=green!15, right=of train] (base)
  {Baseline\\학습\\{\scriptsize XGBoost}};

% ─ Phase 3: TTA ─
\node[box, fill=orange!15, below=1.2cm of caps] (distbuild)
  {테스트 분포\\생성\\{\scriptsize natural/moderate/}\\{\scriptsize attack\_parity}};

\node[box, fill=orange!15, right=of distbuild] (tta)
  {TTA 가중치\\학습\\{\scriptsize Adam, JSD loss}};

\node[box, fill=orange!15, right=of tta] (pred)
  {예측\\{\scriptsize $\sum_e w_e \cdot \ell_e(x)$}\\{\scriptsize argmax}};

% ─ Phase 4: 평가 ─
\node[box, fill=red!10, right=of pred] (eval)
  {평가 \& 시각화\\{\scriptsize A/B/C 플롯}\\{\scriptsize per-class CSV}};

% ─ Oracle ─
\node[box, fill=purple!10, below=0.4cm of train] (oracle)
  {Oracle 예측\\{\scriptsize 상한선}};

% ─ 화살표 ─
\draw[arr] (data)  -- (split);
\draw[arr] (split) -- (caps);
\draw[arr] (caps)  -- (train);
\draw[arr] (train) -- (base);

\draw[arr] (split.south) -- ++(0,-0.35) -| (distbuild.north);
\draw[arr] (distbuild)   -- (tta);
\draw[arr] (tta)         -- (pred);
\draw[arr] (pred)        -- (eval);
\draw[arr] (base.south)  |- (eval.north);

\draw[arr] (train.south) -- (oracle.north);
\draw[arr] (oracle.east) -| (eval.south);

% ─ Phase 레이블 ─
\node[phase, above=0.05cm of data, xshift=0.6cm]
  {\textcolor{blue!60}{Phase 1: 데이터}};
\node[phase, above=0.05cm of caps, xshift=0.8cm]
  {\textcolor{green!50!black}{Phase 2: 학습}};
\node[phase, left=0.1cm of distbuild, yshift=0.1cm]
  {\textcolor{orange!70!black}{Phase 3: TTA}};

\end{tikzpicture}
\end{center}
\end{frame}

% ============================================================
\section{Inclusive Expert 설계}
% ============================================================

\begin{frame}{Inclusive Expert: Cap 기반 설계}

\begin{columns}[T]
\column{0.52\textwidth}
\textbf{설계 원리}
\begin{itemize}
  \item 모든 Expert가 \textbf{전체 15개 클래스}를 분류
  \item Expert $i$의 학습 데이터: 각 클래스 샘플 수를 \\
        $\text{cap}_i$로 제한
  \item $\text{cap}_i$ = 학습셋 클래스 수를 내림차순 정렬 후\\
        $\lceil n_\text{cls} \cdot i / k \rceil$ 번째 클래스의 샘플 수
  \item Expert 0: cap 없음 (자연 불균형 유지)
  \item Expert $k{-}1$: 가장 균등한 분포
\end{itemize}

\vspace{0.4em}
\textbf{Val set도 동일 cap 적용}\\
{\small → early stopping이 해당 expert의 실제 학습 분포에 맞게 동작}

\vspace{0.4em}
\textbf{CIC-IDS2017 예시 ($k=4$)}\\
\begin{tabular}{ccc}
\toprule
Expert & Cap & 학습 데이터 성격 \\
\midrule
E0 & 없음    & 자연 불균형 (benign 454k) \\
E1 & $\sim$25k & 상위 클래스만 균등화 \\
E2 & $\sim$1k  & 대부분 균등 \\
E3 & $\sim$130 & 거의 완전 균등 \\
\bottomrule
\end{tabular}

\column{0.46\textwidth}
\begin{tikzpicture}[scale=0.78]
% 막대 그래프: 클래스별 샘플 수 vs cap
\begin{axis}[
  ybar, bar width=5pt,
  width=7.5cm, height=5.5cm,
  xlabel={\small 클래스 (지지도 내림차순)},
  ylabel={\small 샘플 수 (log)},
  ymode=log,
  ymin=1, ymax=600000,
  xtick={1,2,3,4,5,6,7},
  xticklabels={benign, dos-hulk, portscan, ddos, ftp-pat, ssh-pat, \ldots},
  xticklabel style={font=\tiny, rotate=35, anchor=east},
  ytick={10,100,1000,10000,100000},
  yticklabels={10,100,1k,10k,100k},
  legend style={font=\tiny, at={(0.98,0.98)}, anchor=north east},
  grid=major, grid style={dashed, gray!40},
  title={\small 클래스별 샘플 수와 Cap 수준},
  title style={font=\small},
]
% 실제 샘플 수
\addplot[fill=blue!40] coordinates {
  (1,454537)(2,46025)(3,31786)(4,25606)(5,1588)(6,1179)(7,1100)};
% E1 cap
\addplot[fill=orange!60, opacity=0.7] coordinates {
  (1,25000)(2,25000)(3,25000)(4,25000)(5,1588)(6,1179)(7,1100)};
% E3 cap
\addplot[fill=green!60, opacity=0.7] coordinates {
  (1,130)(2,130)(3,130)(4,130)(5,130)(6,130)(7,130)};
\legend{자연 분포, E1 cap ($\sim$25k), E3 cap ($\sim$130)}
\end{axis}
\end{tikzpicture}
\end{columns}
\end{frame}

% ============================================================
\section{테스트 분포 구성}
% ============================================================

\begin{frame}{테스트 분포 구성: Benign 수 조절}

\begin{columns}[T]
\column{0.45\textwidth}
\textbf{설계 배경}
\begin{itemize}
  \item 공격 클래스 샘플은 고정 (heartbleed 2개 등)
  \item Benign만 다운샘플링으로 조절 가능
  \item Anchor = 두 번째 최다 클래스 (dos-hulk, $\sim$9k in test)
\end{itemize}

\vspace{0.5em}
\begin{tabular}{lrr}
\toprule
분포 & Benign 수 & Benign:Anchor \\
\midrule
\texttt{natural}       & $\sim$91k & 10:1 \\
\texttt{moderate}      & $\sim$27k & 3:1  \\
\texttt{attack\_parity}& $\sim$9k  & 1:1  \\
\bottomrule
\end{tabular}

\vspace{0.5em}
{\small TTA 가중치 학습은 각 분포로,\\
\textbf{최종 평가는 항상 natural 전체 test set}}

\column{0.53\textwidth}
\begin{tikzpicture}[scale=0.82]
\begin{axis}[
  xbar, bar width=9pt,
  width=8cm, height=5.8cm,
  xlabel={\small 샘플 수},
  symbolic y coords={attack\_parity, moderate, natural},
  ytick=data,
  yticklabel style={font=\small},
  xmin=0, xmax=115000,
  xtick={0,20000,40000,60000,80000,100000},
  xticklabels={0,20k,40k,60k,80k,100k},
  legend style={font=\scriptsize, at={(0.98,0.05)}, anchor=south east},
  grid=major, grid style={dashed, gray!40},
  title={\small 분포별 클래스 구성},
  title style={font=\small},
]
\addplot[fill=benignblue!70] coordinates {
  (91000,natural)(27000,moderate)(9000,attack\_parity)};
\addplot[fill=attred!70] coordinates {
  (22000,natural)(22000,moderate)(22000,attack\_parity)};
\legend{Benign, 공격 클래스 (고정)}
\end{axis}
\end{tikzpicture}
\end{columns}

\vspace{0.3em}
\small
\textbf{기대 효과}: natural → E0(unbalanced) 선호 / attack\_parity → E3(balanced) 선호\\
이 패턴이 관찰되면 \textbf{분포 변화에 따라 가중치가 의미 있게 달라진다} = TTA 설계 작동 검증
\end{frame}

% ============================================================
\section{TTA 가중치 학습}
% ============================================================

\begin{frame}{TTA 가중치 학습: MATI/SADE 하이브리드}

\begin{columns}[T]
\column{0.48\textwidth}
\textbf{이론적 배경}

\begin{tabular}{p{1.8cm}p{3.2cm}}
\toprule
SADE (이미지) & cosine similarity 최대화,\newline global weight, 분류 \\
\midrule
MATI (회귀) & MSE$(y^{(1)}, y^{(2)})$ 최소화,\newline global weight, 회귀 \\
\midrule
\textbf{우리 (분류)} & \textbf{JSD$(p^{(1)}, p^{(2)})$ 최소화},\newline \textbf{global weight}, \textbf{logits 집계} \\
\bottomrule
\end{tabular}

\vspace{0.5em}
\textbf{손실 함수}
\[
  \mathcal{L} = \frac{1}{|D_t|}\sum_{x\in D_t}
  \operatorname{JSD}\!\bigl(p^{(1)}_x,\; p^{(2)}_x\bigr)
\]
\[
  p^{(k)}_x = \operatorname{softmax}\!\Bigl(\sum_e \sigma(w_e)\cdot\ell_e(x^{(k)})\Bigr)
\]
\[
  \operatorname{JSD}(p,q)=\tfrac{1}{2}\operatorname{KL}(p\|m)+
  \tfrac{1}{2}\operatorname{KL}(q\|m),\quad m=\tfrac{p+q}{2}
\]

\column{0.50\textwidth}
\begin{center}
\begin{tikzpicture}[
  node distance=0.4cm and 0.5cm,
  box/.style={draw, rounded corners=3pt, minimum width=2cm,
              minimum height=0.55cm, font=\scriptsize, align=center},
  arr/.style={-{Stealth[scale=0.7]}, thick},
]

% 입력 샘플
\node[box, fill=gray!20] (x) {$x$ (mini-batch)};

% 두 view
\node[box, fill=yellow!30, below left=0.7cm and 0.6cm of x] (v0)
  {$x^{(1)}$\\perturbation};
\node[box, fill=yellow!30, below right=0.7cm and 0.6cm of x] (v1)
  {$x^{(2)}$\\perturbation};

\draw[arr] (x) -- (v0);
\draw[arr] (x) -- (v1);

% Expert logits
\node[box, fill=expblue!30, below=0.5cm of v0] (l0)
  {$\ell_e(x^{(1)})$\\{\tiny $E$개 expert logit}};
\node[box, fill=expblue!30, below=0.5cm of v1] (l1)
  {$\ell_e(x^{(2)})$\\{\tiny $E$개 expert logit}};

\draw[arr] (v0) -- (l0);
\draw[arr] (v1) -- (l1);

% Aggregation
\node[box, fill=expgreen!30, below=0.5cm of l0] (a0)
  {$\text{agg}^{(1)}=\sum_e\sigma(w_e)\ell_e^{(1)}$};
\node[box, fill=expgreen!30, below=0.5cm of l1] (a1)
  {$\text{agg}^{(2)}=\sum_e\sigma(w_e)\ell_e^{(2)}$};

\draw[arr] (l0) -- (a0);
\draw[arr] (l1) -- (a1);

% JSD
\node[box, fill=jsdcolor!25,
      below=0.45cm of $(a0)!0.5!(a1)$] (jsd)
  {$\mathcal{L}=\operatorname{JSD}(p^{(1)},p^{(2)})$};

\draw[arr] (a0.south) -- (jsd.north west);
\draw[arr] (a1.south) -- (jsd.north east);

% Adam
\node[box, fill=red!20, below=0.4cm of jsd] (adam)
  {Adam: $w \leftarrow w - \eta\nabla_w\mathcal{L}$\\
   {\tiny 모델 파라미터 frozen}};

\draw[arr] (jsd) -- (adam);

% w 화살표 돌아오기
\draw[arr, dashed, bend left=50] (adam.west)
  to node[left, font=\tiny]{update $w$} (a0.west);

% XGBoost 고정 표시
\node[draw=none, font=\tiny, text=gray,
      right=0.05cm of l0] {(XGBoost, frozen)};

\end{tikzpicture}
\end{center}
\end{columns}
\end{frame}

% ── TTA 알고리즘 상세 ─────────────────────────────────────────
\begin{frame}{TTA 학습 알고리즘 상세}

\begin{columns}[T]
\column{0.50\textwidth}
\textbf{알고리즘 흐름 (\texttt{learn\_aggregation\_weights})}

\begin{itemize}
  \item \textbf{Step 1.} $w \leftarrow \mathbf{0}_E$ 초기화 (softmax 후 균등)
  \item \textbf{Step 2.} For each epoch:
        $D_t$ mini-batch 추출
        $\to$ 두 독립 perturbation $x^{(1)}, x^{(2)}$
        $\to$ 각 Expert logit 추출 (\texttt{output\_margin})
        $\to$ JSD loss 계산 \& Adam step
  \item \textbf{Step 3.} Early stop: $\min_e w_e < 0.05$
  \item \textbf{Step 4.} 학습된 $w$를 natural test set에 적용
\end{itemize}

\vspace{0.3em}
\textbf{역전파 경로}\\
$\mathcal{L} \to p^{(1,2)} \to \text{agg} \xrightarrow{\partial/\partial w} w$\\
{\small XGBoost logit은 \textbf{상수 텐서}로 취급\\
→ XGBoost 역전파 불필요}

\column{0.48\textwidth}
\textbf{예측 (\texttt{predict\_with\_weights})}

\[
  \hat{y} = \arg\max_c \sum_{e=1}^{E} w_e \cdot \ell_e(x)_c
\]

softmax 이전 logit을 가중합 → argmax\\
(확률을 가중합하는 것보다 이론적으로 타당)

\vspace{0.6em}
\textbf{하이퍼파라미터 (CLI)}

\begin{tabular}{ll}
\toprule
\texttt{--tta\_epochs} & 5 \\
\texttt{--tta\_lr}     & 0.01 \\
\texttt{--tta\_batch}  & 256 \\
\texttt{--num\_experts}& 4 \\
\texttt{--noise\_std}  & 0.1 \\
\bottomrule
\end{tabular}

\vspace{0.5em}
\textbf{MATI vs 우리 설정 차이점}\\
\begin{itemize}\small
  \item MATI: regression 스칼라 $\to$ MSE
  \item 우리: 분류 확률 벡터 $\to$ JSD\\
        (대칭, bounded, 정보이론적 근거)
  \item MATI: GMM region별 전문화\\
  \item 우리: cap 수준별 전문화
\end{itemize}
\end{columns}
\end{frame}

% ============================================================
\section{평가 및 출력}
% ============================================================

\begin{frame}{평가 파이프라인 및 출력물}

\begin{columns}[T]
\column{0.5\textwidth}
\textbf{평가 구조}

\begin{tikzpicture}[
  node distance=0.3cm and 0.4cm,
  box/.style={draw, rounded corners=3pt, minimum width=3.4cm,
              minimum height=0.6cm, font=\scriptsize, align=center,
              fill=lightgray},
  arr/.style={-{Stealth[scale=0.7]}, thick},
]

\node[box, fill=blue!10] (te)
  {Natural test set ($X_\text{te}$, $y_\text{te}$)};

\node[box, fill=green!10, below=0.3cm of te] (base)
  {Baseline 예측 (\texttt{baseline.predict})};

\node[box, fill=purple!10, below=0.3cm of base] (orc)
  {Oracle 예측 (true class 최고확률 expert)};

\node[box, fill=orange!10, below=0.5cm of orc] (tta_nat)
  {TTA-natural 예측};
\node[box, fill=orange!10, below=0.3cm of tta_nat] (tta_mod)
  {TTA-moderate 예측};
\node[box, fill=orange!10, below=0.3cm of tta_mod] (tta_att)
  {TTA-attack\_parity 예측};

\draw[arr] (te)   -- (base);
\draw[arr] (te)   -- (orc);
\draw[arr] (te)   -- (tta_nat);
\draw[arr] (te)   -- (tta_mod);
\draw[arr] (te)   -- (tta_att);

\node[box, fill=red!10, right=0.5cm of orc] (eval)
  {비교 평가\\(Macro/Weighted F1)};

\draw[arr] (base)    -- (eval);
\draw[arr] (orc)     -- (eval);
\draw[arr] (tta_nat) -- (eval);
\draw[arr] (tta_mod) -- (eval);
\draw[arr] (tta_att) -- (eval);

\end{tikzpicture}

\column{0.48\textwidth}
\textbf{출력 파일}

\begin{tabular}{ll}
\toprule
파일 & 내용 \\
\midrule
\texttt{A\_f1\_comparison.png} & 클래스별 F1 비교 막대 \\
\texttt{A\_f1\_\{dist\}.csv}   & 분포별 클래스 F1 수치 \\
\texttt{B\_expert\_behavior.png} & Expert 예측 테이블\\
                                 & + 분포별 가중치 열 \\
\texttt{B\_expert\_behavior.csv} & 위 수치 \\
\texttt{C\_global\_weights.png}  & 분포별 학습 가중치\\
                                 & 막대 그래프 \\
\texttt{C\_global\_weights.csv}  & 위 수치 \\
\texttt{baseline\_vs\_moe}       & Baseline vs TTA \\
\texttt{\_per\_class\_\{dist\}} & 클래스별 정밀/재현/F1 \\
\texttt{.png/.csv}               & 컬러 표 + CSV \\
\bottomrule
\end{tabular}

\vspace{0.3em}
{\small 분포별로 별도 결과 파일 생성 (\texttt{\_natural},\\
\texttt{\_moderate}, \texttt{\_attack\_parity} suffix)}
\end{columns}
\end{frame}

% ============================================================
\section{현재 한계 및 관찰}
% ============================================================

\begin{frame}{현재 한계 및 관찰 (실험 결과 기준)}

\begin{columns}[T]
\column{0.5\textwidth}
\textbf{관찰된 문제}

\begin{itemize}
  \item \textbf{Early stop이 epoch 1에서 발생}\\
        E0 가중치 $\approx 0.9$ → 나머지 $< 0.05$\\
        원인: \texttt{tta\_lr=0.01}이 너무 큼

  \item \textbf{E0가 모든 분포에서 지배적}\\
        E0 = 가장 많은 데이터 → logit이 perturbation에\\
        가장 안정적 → JSD 최소화 = E0 선택\\
        {\small(올바른 최적화지만 원하는 결과 아님)}

  \item \textbf{Tail 클래스 미개선}\\
        bot, web-attack-xss, sql-injection\\
        Expert 예측 자체가 benign으로 수렴
\end{itemize}

\column{0.48\textwidth}
\textbf{근본 원인 분석}

\begin{tabular}{p{2cm}p{3.5cm}}
\toprule
문제 & 원인 \\
\midrule
Epoch 1 stop & lr=0.01 → 첫 배치에서\\
             & weight collapse \\
\midrule
E0 지배      & Generalist expert는\\
             & 항상 안정적 → 분포에\\
             & 무관하게 선택됨 \\
\midrule
Tail 미개선  & Expert 설계 문제:\\
             & 모든 expert가\\
             & benign으로 수렴 \\
\bottomrule
\end{tabular}

\vspace{0.4em}
\textbf{제안 수정}\\
{\small
• \texttt{tta\_lr}: 0.01 → 0.0001\\
• Early stop threshold: 0.05 → 0.01\\
• Entropy 정규화 추가:\\
  $\mathcal{L} = \text{JSD} - \lambda H(\mathbf{w})$
}
\end{columns}

\vspace{0.3em}
\begin{alertblock}{\small 핵심 한계}
\small TTA는 Expert 품질이 서로 달라야 작동함.
Expert들이 모두 benign으로 수렴하면 어떤 routing도 의미 없음.
\textbf{선결 과제: Expert 학습 자체의 tail-class 성능 개선.}
\end{alertblock}
\end{frame}

% ── 요약 ─────────────────────────────────────────────────────
\begin{frame}{요약}

\begin{block}{설계 요소}
\begin{itemize}
  \item \colorbox{blue!15}{(1)} \textbf{Inclusive Expert}: $k$개 전문가, 각자 다른 cap으로 균형 수준 차별화
  \item \colorbox{green!20}{(2)} \textbf{Val set cap 적용}: early stopping이 학습 분포 기준으로 동작
  \item \colorbox{orange!20}{(3)} \textbf{테스트 분포 3종}: natural / moderate / attack\_parity (benign 수 조절)
  \item \colorbox{purple!15}{(4)} \textbf{TTA}: JSD 최소화로 global $\mathbf{w}[E]$ 학습 (MATI 원리 + 분류용 JSD)
  \item \colorbox{red!15}{(5)} \textbf{예측}: logit 가중합 후 argmax (softmax 이전 집계)
\end{itemize}
\end{block}

\vspace{0.5em}
\begin{tabular}{lll}
\toprule
방법  & 핵심 설계 & 우리와의 관계 \\
\midrule
MATI  & 다중 분포 + MSE 일관성, 회귀  & 다중 분포 설계 차용 \\
SADE  & global 스칼라 weight, 분류  & 가중치 학습 방식 차용 \\
\textbf{우리} & \textbf{다중 분포 + JSD 일관성, 분류} & MATI + SADE 혼합 \\
\bottomrule
\end{tabular}

\end{frame}

\end{document}
